\section{Reference Counting}

% GOAL: This introduction should say, we will:
%  - explain reference counting is
%  - explain how "Counting immutable beans" work
%  - then finally list the constraints required by the "counting immutable beans" algorithm

This section introduces reference counting as a memory management strategy and explains how the approach of \citeauthor*{ullrich_counting_2021}~\cite{ullrich_counting_2021} uses \textbf{reset}/\textbf{reuse} instructions to reduce memory allocations and reference ownership information to reduce the overhead of reference counting.
We then present the constraints required by the algorithm of \citeauthor*{ullrich_counting_2021}~\cite{ullrich_counting_2021}, since these become important when we later adapt the algorithm to Rizzo in Section~\ref{section:refcount}.

% In Rizzo signal combinators create new signals as an intermediate step in the computation. These intermediate signals are often only used for a short period of time, to copy over the new computed value to the head of the original signal. Given the intermediate nature of these signals, it is important to reclaim the memory used by them as soon as they are no longer needed, otherwise they will also be triggered alongside the original signal. The retrigger of the signals will lead to performance issues, memory leaks and unexpected side effects \todo{We will get back to why this is important in section XXX where we describe the structure of the heap}. To address the issue, we need a way to efficiently manage the memory used by these intermediate signals and ensure that it is reclaimed as soon as it is no longer needed.

% One way to achieve this is through reference counting, which is a technique for managing memory in programming languages. In this section, we will discuss the reference counting technique proposed by \citetitle{ullrich_counting_2021} \cite{ullrich_counting_2021} and how we can adapt it to Rizzo to efficiently manage the memory used by intermediate signals.

\subsection{What is reference counting}
% what does it mean to reference count
Early work on reference counting like that of \citeauthor*{collins_method_1960} \cite{collins_method_1960}
%and \citetitle{bacon_computing_2001} \cite{bacon_computing_2001} 
introduced the idea of reference counting as an alternative to other memory management methods like tracing garbage collection. The basic idea of reference counting is to attach a counter to each object in memory that keeps track of how many references point to the object. When a new reference is created, the counter is incremented, and when a reference is destroyed, the counter is decremented. When the counter reaches zero, it means that there are no more references to the object, and it can be safely deallocated.

Early implementations were implemented in the runtime system of the programming language. This was a common approach, but it is hard to optimize.
Newer techniques, like the one proposed by \citeauthor*{ullrich_counting_2021}~\cite{ullrich_counting_2021}, do static analysis during compilation in order to achieve better performance and more efficient memory management. 

By using new \textbf{reset} and \textbf{reuse} instructions and ownership annotations, the compiler can determine when memory of immutable data can be safely overwritten, and when increment/decrement instructions can or cannot be safely omitted.

% which parts of the program are responsible for managing the memory of each object, the compiler can determine when an object can be safely deallocated, and optimize reference counting operations accordingly. The optimisation comes from knowing when work (incrementing or decrementing reference counts) can be safely omitted, and when it cannot.

Later in this section, we will be using the term \textit{RC token} (Reference Counting token).
An RC token is one counted reference to an object.
A token is created when the reference count is incremented, and it is consumed when the reference count is decremented, either directly or by being passed to a context that eventually decrements it, such as a function or a memory allocated value.

% of a reference to an object. We say that when a new reference is created, an RC token is created and when a reference is destroyed, an RC token is consumed.

% circular references? bacon papers solution to this? we don't have this problem in Rizzo because of mutable heads, but maybe we can still talk about it?

\subsubsection{Naïve reference counting}

In a naïve reference counting implementation, every time a new reference to an object is created, the reference count of that object is incremented. This means the creator of the object has to increment the reference count, and every time the object is taken as input, the reference count has to be incremented. We have to do this because we cannot be sure if the object is shared or not, and we have to be safe. This can lead to a lot of unnecessary increments and decrements, especially in cases where many short-lived objects are created and destroyed, or an object is passed multiple times to a function.

Using a notation close to the reference counted intermediate language generated by the compiler, we can write an example of a naïve reference counting implementation. Let us assume we have a function $f$ that takes an object as input, and then we use it from the main function.

\begin{spacing}{0.8}
\begin{Verbatim}
fun f x =
    inc x;             -- ref: 2
    ...                -- use x
    dec x;             -- ref: 1

let main =
    let z = Obj() in   -- start count 0
    inc z;             -- ref: 1
    f z;
    dec z;             -- end count 0
\end{Verbatim}
\end{spacing}

As we can see, the reference to the object is being incremented twice and decremented twice. We increment it before we pass it to $f$, and then have to increment it again inside $f$ because we don't know how the object is used. We then have to decrement it twice, once inside $f$ before we return, and once after we call $f$ in the main function. This is quite inefficient, especially if we passed $z$ further to another function inside $f$.


\subsection{Counting immutable beans}

\citeauthor*{ullrich_counting_2021} \cite{ullrich_counting_2021} gain their performance improvements by using a combination of ownership annotations, and new \textbf{reset} and \textbf{reuse} instructions.
The ownership annotations in \citetitle{ullrich_counting_2021} \cite{ullrich_counting_2021} work by distinguishing between $\mathbb{O}wned$ and $\mathbb{B}orrowed$ references, and let the compiler track which parts of the program are responsible for managing the memory of each object. The compiler can then use this knowledge to determine when an object needs to be incremented or decremented, which reduces the number of reference counting operations that need to be performed.
For owned references we need to ensure that the object is alive as long as we need it, which is similar to the naïve reference counting implementation.
If we instead borrow an object, then we assume that the caller has the responsibility to ensure that the object is alive as long as we need it. This means we can avoid incrementing the reference count of the object, which is a great performance improvement in nested scopes where all the functions can borrow their parameters. In this case, we can increment the reference count once at the top of the scope, instead of incrementing it for each function call as we would need to if it was owned.
We can typically borrow an object when we only need to read from it, and then own the object if we intend to reuse it.

We can then extend this with the new \textbf{reset} and \textbf{reuse} instructions. These instructions allow the compiler to optimize memory allocations of new objects by reusing memory from old objects that are no longer needed, instead of always allocating new memory for each object. This leads to performance improvements, especially in cases where many short-lived objects or nested objects are created and destroyed.

The paper defines three main steps for implementing their reference counting technique, where each step depends on the previous one:

\begin{itemize}
    \item Inserting \textbf{reset}/\textbf{reuse} pairs
    \item Inferring borrowed parameters
    \item Inserting \textbf{inc}/\textbf{dec} instructions
\end{itemize}

% (1) Inserting reset/reuse pairs (Section 5.1)  
% (2) Inferring borrowed parameters (Section 5.2)  
% (3) Inserting inc/dec instructions (Section 5.3)


\subsubsection{Reference Counting with Ownership}\label{section:beans:examples}

To show how ownership annotations and \textbf{reset}/\textbf{reuse} instructions can optimize reference counting, we will go through some examples. These examples will use similar notation as the one we used in the naïve example. The notation is based on our intermediate representation, where \textbf{inc} and \textbf{dec} are used to increment and decrement the reference count of a variable, while \textbf{reset} and \textbf{reuse} are used to reclaim and recycle storage.
% Using ownership annotations we can reduce the number of increments and decrements needed, by allowing the compiler to determine when an object can be safely deallocated, and to optimize reference counting operations accordingly.
% % The following examples use a notation close to the reference counted intermediate language generated by the compiler. 
% The instructions \textbf{inc} and \textbf{dec}, increment and decrement the reference count of a variable, while \textbf{reset} and \textbf{reuse} are used to reclaim and recycle storage.

% Lets revisit the example of the function $f$ from the naïve approach, We will now show how ownership annotations can optimize the reference counting in this example. We first show an example where the parameter is owned, and then we show how we can further optimize it by borrowing the parameter instead.

% With our ownership annotations, all objects start with a reference count of 1, which implicitly means there is an \textbf{inc} instruction.

% \begin{spacing}{0.8}
% \begin{Verbatim}
% x: owned
% fun f x =
%     ...                -- use x
%     dec x;             -- ref: 0

% let main =
%     let z = Obj() in   -- start count 1
%     f z;
% \end{Verbatim}
% \end{spacing}




As a first example, consider the expression below where the same variable is passed twice to a function:

\begin{verbatim}
let z = f y y in
ret z
\end{verbatim}

Assume that the function $f$ takes two parameters and consumes an RC token for each of them. Consuming an RC token for a parameter will mark the parameter as $\mathbb{O}wned$, which means that the caller is responsible for the reference count of the objects. In this case, since $y$ is passed twice to $f$, both occurrences of $y$ are treated as owned, and both will consume an RC token. The compiler must then insert an \textbf{inc} before the call to make the second owned use valid:

\begin{verbatim}
let y = Obj() in
inc y;
let z = f y y in
ret z
\end{verbatim}

No matching \textbf{dec} is needed afterwards, since both RC tokens are consumed by the call itself. This is the safe but cautious behaviour we get when both arguments are treated as owned.

Borrowing removes some of this overhead. If the first parameter of $f$ is borrowed and the second is owned, then only one of the two uses of $y$ consumes an RC token. The borrowed use still requires the object to stay alive for the duration of the call, so the compiler inserts a temporary increment before the call and a matching decrement afterwards:

\begin{verbatim}
let y = Obj() in
inc y;
let z = f y y in
dec y;
ret z
\end{verbatim}

Compared to the previous example, this makes the ownership discipline more precise: one occurrence of $y$ is only observed, while the other is actually consumed. The compiler can therefore avoid treating both arguments as fully owned.

The more interesting optimisation is enabled by the combination of \textbf{reset} and \textbf{reuse}. In Rizzo this matters when a combinator stops using an input signal and immediately constructs a replacement signal. Consider the function below:

\begin{lstlisting}
fun stop f (x :: xs) : ('a -> Bool) -> Signal 'a -> Signal 'a =
    if f x 
    then x :: never
    else x :: (stop f |> xs)
\end{lstlisting}

In the `then' branch, the input signal $(x :: xs)$ is not used. This makes it possible to reuse the memory cell to store the new signal $x :: \mathsf{never}$, avoiding an allocation. 

Not every signal combinator admits this optimisation. For a function such as $map$, the input signal must remain alive because its tail is needed to produce future values:

\begin{lstlisting}
fun map f (x :: xs) = 
    f x :: (map f |> xs)
\end{lstlisting}

Here the delayed computation still depends on $xs$, so reusing the storage of $s$ would destroy data that is needed later. This illustrates why reuse opportunities for signals in Rizzo are limited: they mainly arise in combinators that terminate, switch, or otherwise stop forwarding the original signal.



\subsection{AST structure requirements}\label{sec:background:beans:assumptions}
% transformations that we need

% (1) Inserting reset/reuse pairs (Section 5.1)  
% (2) Inferring borrowed parameters (Section 5.2)  
% (3) Inserting inc/dec instructions (Section 5.3)

% We assume that
% • all constructor applications are fully applied by eta-expanding them. ?? not needed till we add custom types and constructors?
% • no constant applications are over-applied by splitting them into two applications where necessary.
% • all variable applications take only one argument, again by splitting them where necessary. While this simplification can introduce additional allocations of intermediary partial applications, it greatly simplifies the presentation of our operational semantics. All presented program transformations can be readily extended to a system with n-ary variable applications, which are handled analogously to n-ary constant applications.
% • every function abstraction has been lambda-lifted to a toplevel constant c.
% • trivial bindings let x = y have been eliminated through copy propagation.
% • all dead let bindings have been removed.
% • all parameter and let names of a function are mutually distinct. Thus we do not have to worry about name capture.

% how we are going to use it and what we use and don't use from the paper
% for example we don't (need to) fully apply constructors? move this somewhere later?

To implement the reference counting technique proposed by \citeauthor*{ullrich_counting_2021}~\cite{ullrich_counting_2021}, the shape of the abstract syntax tree (AST) must satisfy the following properties:

\begin{itemize}
    \item Eta-expanding constructor applications to ensure that all constructors are fully applied.
    \item Splitting over-applied constant applications into multiple applications to ensure that all constant applications are fully applied.
    \item Splitting variable applications into multiple applications to ensure that all variable applications take only one argument.
    \item Lambda-lifting all function abstractions to the top level to ensure that all functions are defined at the top level of the program.
    \item Eliminating trivial bindings of the form \lstinline|let x = y| through copy propagation, to simplify the AST and reduce unnecessary bindings.
    \item Removing all dead let bindings to clean up the AST and remove unused variables.
    \item Ensuring that all parameter and let names of a function are mutually distinct to avoid name capture issues during transformations.
\end{itemize}

We define two separate abstract syntaxes: one for the Rizzo language, described in Section~\ref{section:rizzolang}, and another for the intermediate representation (IR) used in the insertion of reference counting instructions, described in Section~\ref{section:refcount}.
To ensure the properties above and to convert from the Rizzo AST to the IR, we define a series of transformations in Section~\ref{section:refcount:transforms}.
